% ---
% id: receiver/lem:exactwindow
% kind: lem
% title: "Exact window with an analytic remainder"
% summary: >-
%   Lower estimates of the exact block bound on finitely many blocks above the
%   cutoff, combined with an admissible remainder for the blocks beyond them,
%   give a weight in [0,1] that bounds the multiclick element on every block
%   above the cutoff. Three admissible choices are named.
% prominence: standard
% status: implemented
% uses: [thm:multiclick-bound, lem:exactbasic, lem:windowtail, lem:collapse]
% assumes: []
% refs: []
% provenance: newly-derived
% implements: certiqs_key.receiver.bound:exact_composite
% lean:
% ---

The lemma turns information about finitely many blocks into a bound uniform
over all blocks above the cutoff: a lower estimate on each exact block in a
window, and a remainder that covers every block beyond it. The three choices
listed at its end are all this contract uses.

\begin{lemma}\label{lem:exactwindow}
Let the hypotheses of Theorem~\ref{thm:multiclick-bound} hold together with
\eqref{eq:mcpos}, let $N\ge0$, and let
$m^{\ast}>N$ be an integer. Let $\ell(n)$, $N<n\le m^{\ast}$, be reals with
\begin{equation}\label{eq:ellhyp}
  \ell(n)\ \le\ \Bnd^{\ast}(n)\qquad\text{for }N<n\le m^{\ast},
\end{equation}
and let $\tau$ be a real with
\begin{equation}\label{eq:tailadm}
  \tau\le0
  \qquad\text{or}\qquad
  \tau\le\inf_{n>m^{\ast}}\Bnd^{\ast}(n).
\end{equation}
Then
\begin{equation}\label{eq:exactcomposite}
  P\ \coloneqq\ \max\Big\{0,\ \min\Big\{\ \min_{N<n\le m^{\ast}}\ell(n),
   \ \ \tau\ \Big\}\Big\}
\end{equation}
satisfies $0\le P\le1$ and
\begin{equation}\label{eq:exactstatement}
  \blk{M}{\gtN}_{\mc}\ \succeq\ P\,\id^{(>N)} .
\end{equation}
Three choices are admissible:
\begin{enumerate}[label=(\roman*),leftmargin=3.0em]
\item $\tau=T_{m^{\ast}}$ of \eqref{eq:windowtail}, formed from the bound
  data $\big(\gamma,\{\beta_k\},c,\{d_k\}\big)$ that
  Theorem~\ref{thm:multiclick-bound} forms by \eqref{eq:cd}, satisfies
  \eqref{eq:tailadm} with no further hypothesis;
\item if the bound data satisfy \eqref{eq:collapsehyp}, then
  $\tau=\Bnd(m^{\ast}+1)$ satisfies \eqref{eq:tailadm};
\item $\ell(n)=\Bnd(n)$ satisfies \eqref{eq:ellhyp}, and the pointwise maximum
  of finitely many admissible choices of $\ell$ is admissible.
\end{enumerate}
\end{lemma}

\begin{proof}
\emph{Step 1 (range of $P$).} $0\le P$ is the outer maximum. For $P\le1$:
either $P=0$, or $P\le\ell(N+1)\le\Bnd^{\ast}(N+1)\le1$ by \eqref{eq:ellhyp}
and \eqref{eq:exactbasic}, which needs no more than the hypotheses of
Theorem~\ref{thm:multiclick-bound}, the window being non-empty because
$m^{\ast}>N$.

\emph{Step 2 (the bound \eqref{eq:exactstatement}).} If $P=0$ the assertion is
$\blk{M}{\gtN}_{\mc}\succeq0$, and this is the one step that uses
\eqref{eq:mcpos}: that hypothesis gives $\blk{M}{\eqn}_{\mc}\succeq0$ on every
block (Lemma~\ref{lem:exactbasic}), hence on their direct sum. So let $P>0$.
Then
$P=\min\{\min_{N<n\le m^{\ast}}\ell(n),\ \tau\}$, and in particular
$\tau\ge P>0$, so the first disjunct of \eqref{eq:tailadm} fails and the second
holds. Take any $n>N$. If $n\le m^{\ast}$ then $\Bnd^{\ast}(n)\ge\ell(n)\ge P$
by \eqref{eq:ellhyp}; if $n>m^{\ast}$ then
$\Bnd^{\ast}(n)\ge\inf_{j>m^{\ast}}\Bnd^{\ast}(j)\ge\tau\ge P$. Either way
$\lmin(\blk{M}{\eqn}_{\mc})\ge P$, that is
$\blk{M}{\eqn}_{\mc}\succeq P\,\id^{(=n)}$; the blocks with $n>N$ are
orthogonal direct summands, and a lower bound holding on each holds on the
direct sum.

\emph{Step 3 (choice (i)).} For $n>m^{\ast}$, $d_k\le1$ gives
$d_k^{\,n}\le d_k^{\,m^{\ast}}$ and $\gamma c^{\,n}\ge0$, so
$\Bnd(n)\ge1-\sum_k\beta_kd_k^{\,n}\ge
1-\sum_k\beta_kd_k^{\,m^{\ast}}=T_{m^{\ast}}$, the tail step of
Lemma~\ref{lem:windowtail}\,(i); and $\Bnd^{\ast}(n)\ge\Bnd(n)$ by
Lemma~\ref{lem:exactbasic}. Hence
$\inf_{n>m^{\ast}}\Bnd^{\ast}(n)\ge T_{m^{\ast}}$.

\emph{Step 4 (choice (ii)).} If $\Bnd(m^{\ast}+1)\le0$, the first disjunct of
\eqref{eq:tailadm} holds. Otherwise $\Bnd(m^{\ast}+1)>0$ and
Lemma~\ref{lem:collapse}, applied with cutoff $m^{\ast}$ in place of $N$,
gives $\inf_{n>m^{\ast}}\Bnd(n)=\Bnd(m^{\ast}+1)$; with
$\Bnd^{\ast}\ge\Bnd$ pointwise this is
$\inf_{n>m^{\ast}}\Bnd^{\ast}(n)\ge\Bnd(m^{\ast}+1)$, the second disjunct.

\emph{Step 5 (choice (iii)).} The first statement is
Lemma~\ref{lem:exactbasic}. For the second, a maximum of finitely many
numbers each at most $\Bnd^{\ast}(n)$ is at most $\Bnd^{\ast}(n)$.
\end{proof}
